\documentclass[10pt,a4paper]{amsart}
\usepackage{amsmath,amssymb,amsthm,graphicx}
\usepackage{verbatim}
\newenvironment{reference}{\comment}{\endcomment}
%\newenvironment{reference}{}{}
\newenvironment{slogan}{\comment}{\endcomment}
\newenvironment{history}{\comment}{\endcomment}

% For commutative diagrams we use Xy-pic
\usepackage[all]{xy}

% We use 2cell for 2-commutative diagrams.
\xyoption{2cell}
% Theorem environments.
%
\theoremstyle{plain}
\newtheorem{theorem}[subsection]{Theorem}
\newtheorem{proposition}[subsection]{Proposition}
\newtheorem{lemma}[subsection]{Lemma}

\theoremstyle{definition}
\newtheorem{definition}[subsection]{Definition}
\newtheorem{example}[subsection]{Example}
\newtheorem{exercise}[subsection]{Exercise}
\newtheorem{situation}[subsection]{Situation}

\theoremstyle{remark}
\newtheorem{remark}[subsection]{Remark}
\newtheorem{remarks}[subsection]{Remarks}

\numberwithin{equation}{subsection}

% Macros
%
\def\lim{\mathop{\mathrm{lim}}\nolimits}
\def\colim{\mathop{\mathrm{colim}}\nolimits}
\def\Spec{\mathop{\mathrm{Spec}}}
\def\Hom{\mathop{\mathrm{Hom}}\nolimits}
\def\Ext{\mathop{\mathrm{Ext}}\nolimits}
\def\SheafHom{\mathop{\mathcal{H}\!\mathit{om}}\nolimits}
\def\SheafExt{\mathop{\mathcal{E}\!\mathit{xt}}\nolimits}
\def\Sch{\mathit{Sch}}
\def\Mor{\mathop{\mathrm{Mor}}\nolimits}
\def\Ob{\mathop{\mathrm{Ob}}\nolimits}
\def\Sh{\mathop{\mathit{Sh}}\nolimits}
\def\NL{\mathop{N\!L}\nolimits}
\def\CH{\mathop{\mathrm{CH}}\nolimits}
\def\proetale{{pro\text{-}\acute{e}tale}}
\def\etale{{\acute{e}tale}}
\def\QCoh{\mathit{QCoh}}
\def\Ker{\mathop{\mathrm{Ker}}}
\def\Im{\mathop{\mathrm{Im}}}
\def\Coker{\mathop{\mathrm{Coker}}}
\def\Coim{\mathop{\mathrm{Coim}}}

% Boxtimes
%
\DeclareMathSymbol{\boxtimes}{\mathbin}{AMSa}{"02}

%
% Macros for moduli stacks/spaces
%
\def\QCohstack{\mathcal{QC}\!\mathit{oh}}
\def\Cohstack{\mathcal{C}\!\mathit{oh}}
\def\Spacesstack{\mathcal{S}\!\mathit{paces}}
\def\Quotfunctor{\mathrm{Quot}}
\def\Hilbfunctor{\mathrm{Hilb}}
\def\Curvesstack{\mathcal{C}\!\mathit{urves}}
\def\Polarizedstack{\mathcal{P}\!\mathit{olarized}}
\def\Complexesstack{\mathcal{C}\!\mathit{omplexes}}
% \Pic is the operator that assigns to X its picard group, usage \Pic(X)
% \Picardstack_{X/B} denotes the Picard stack of X over B
% \Picardfunctor_{X/B} denotes the Picard functor of X over B
\def\Pic{\mathop{\mathrm{Pic}}\nolimits}
\def\Picardstack{\mathcal{P}\!\mathit{ic}}
\def\Picardfunctor{\mathrm{Pic}}
\def\Deformationcategory{\mathcal{D}\!\mathit{ef}}
\usepackage{cleveref}
\crefname{originaldefinition}{original definition}{original definitions}
\Crefname{originaldefinition}{Original definition}{Original definitions}
\makeatletter
\newlabel{definition-point}{{00Y5}{source}{}{sourcepointer-1}{}}
\newlabel{definition-point@cref}{{[originaldefinition][0][]00Y5}{[1][0][]source}{}{}{}}
\makeatother
\crefname{originallemma}{original lemma}{original lemmas}
\Crefname{originallemma}{Original lemma}{Original lemmas}
\makeatletter
\newlabel{lemma-points-recover}{{00Y6}{source}{}{sourcepointer-2}{}}
\newlabel{lemma-points-recover@cref}{{[originallemma][0][]00Y6}{[1][0][]source}{}{}{}}
\makeatother
\crefname{originallemma}{original lemma}{original lemmas}
\Crefname{originallemma}{Original lemma}{Original lemmas}
\makeatletter
\newlabel{lemma-point-pushforward-sheaf}{{00Y8}{source}{}{sourcepointer-3}{}}
\newlabel{lemma-point-pushforward-sheaf@cref}{{[originallemma][0][]00Y8}{[1][0][]source}{}{}{}}
\makeatother
\crefname{originallemma}{original lemma}{original lemmas}
\Crefname{originallemma}{Original lemma}{Original lemmas}
\makeatletter
\newlabel{lemma-neighbourhoods-directed}{{00YB}{source}{}{sourcepointer-4}{}}
\newlabel{lemma-neighbourhoods-directed@cref}{{[originallemma][0][]00YB}{[1][0][]source}{}{}{}}
\makeatother
\title{Points of a site with finite limits are functors commuting with finite limits and coverings}
\author{The Stacks Project Authors}
\date{stacks-project:sites.tex:local-snapshot-20261012; GNU FDL 1.2 or later}
\begin{document}
\maketitle
\noindent\textit{Editorial note.} All delivered mathematical text is the author's own, copied from the pinned Stacks Project chapter file sites-cohomology.tex; no statement or proof step was rewritten, reordered or removed.\par\medskip
\noindent\textit{Editorial note.} The proof is the author's complete argument; every step the author writes out is inside the counted span and no closing sentence is a placeholder.\par\medskip
\noindent\textit{Editorial note.} Every labelled reference occurring inside the counted spans is carried in assembly.external\_references as a non-load-bearing navigation pointer (same-chapter ones at their definition site, cross-chapter ones at the invocation inside the counted span), so no load-bearing mathematics is replaced by a stub. No bibliography citation occurs inside the counted spans.\par\medskip
\noindent\textit{Editorial note.} The checked chapter list for same-name collisions is the set of stacks chapters staged in this repository: chow, cohomology, divisors, etale-cohomology, flat, more-algebra, more-morphisms, morphisms, spaces-more-morphisms, varieties together with this card's chapter sites-cohomology.\par\medskip
\noindent\textit{Editorial note.} The preamble is the chapter's own retained minimal preamble (sources/stacks-sites-cohomology/retained-preamble.tex), a verbatim excerpt of the native preamble: the Xy-pic setup and the native theorem declarations with the mathematical macros.\par\medskip
\noindent\textit{Allowed standard prerequisites (editorial).}\begin{itemize}
\item The definition of a point of a site. Same chapter, Definition definition-point. Supplies the conditions shown to be equivalent to the two displayed ones.
\item Stalks of representable sheaves. Same chapter, Lemma lemma-points-recover and Lemma lemma-point-pushforward-sheaf. Identify the functor with the composite of the representable functor, sheafification and the stalk functor.
\item Directedness of the category of neighbourhoods. Same chapter, Lemma lemma-neighbourhoods-directed. Gives commutation with finite limits for a functor satisfying the two conditions.
\end{itemize}\medskip
\ifdefined\hypertarget\hypertarget{sourcepointer-1}{}\fi
\noindent\textit{Source reference.} Original definition 00Y5: sources/stacks-sites/pinned-original-source.tex, source lines 8539--8539. Navigation pointer invoked at proof line 8539; the counted span contains the pointer only and the cited result is not reproduced inside it, so the pointer is registered as a non-load-bearing navigation stub in the established form for Stacks material. No bibliography citation is involved.\par\medskip
\ifdefined\hypertarget\hypertarget{sourcepointer-2}{}\fi
\noindent\textit{Source reference.} Original lemma 00Y6: sources/stacks-sites/pinned-original-source.tex, source lines 8541--8541. Navigation pointer invoked at proof line 8541; the counted span contains the pointer only and the cited result is not reproduced inside it, so the pointer is registered as a non-load-bearing navigation stub in the established form for Stacks material. No bibliography citation is involved.\par\medskip
\ifdefined\hypertarget\hypertarget{sourcepointer-3}{}\fi
\noindent\textit{Source reference.} Original lemma 00Y8: sources/stacks-sites/pinned-original-source.tex, source lines 8542--8542. Navigation pointer invoked at proof line 8542; the counted span contains the pointer only and the cited result is not reproduced inside it, so the pointer is registered as a non-load-bearing navigation stub in the established form for Stacks material. No bibliography citation is involved.\par\medskip
\ifdefined\hypertarget\hypertarget{sourcepointer-4}{}\fi
\noindent\textit{Source reference.} Original lemma 00YB: sources/stacks-sites/pinned-original-source.tex, source lines 8559--8559. Navigation pointer invoked at proof line 8559; the counted span contains the pointer only and the cited result is not reproduced inside it, so the pointer is registered as a non-load-bearing navigation stub in the established form for Stacks material. No bibliography citation is involved.\par\medskip
\begin{proposition}
\label{proposition-point-limits}
Let $\mathcal{C}$ be a site. Assume that finite limits exist
in $\mathcal{C}$. (I.e., $\mathcal{C}$ has fibre products, and a
final object.) A point $p$ of such a site $\mathcal{C}$
is given by a functor $u : \mathcal{C} \to \textit{Sets}$ such that
\begin{enumerate}
\item $u$ commutes with finite limits, and
\item if $\{U_i \to U\}$ is a covering, then
$\coprod_i u(U_i) \to u(U)$ is surjective.
\end{enumerate}
\end{proposition}

\begin{proof}
Suppose first that $p$ is a point (Definition \ref{definition-point})
given by a functor $u$. Condition (2) is satisfied directly from
the definition of a point. By Lemma \ref{lemma-points-recover}
we have $(h_U)_p = u(U)$. By Lemma \ref{lemma-point-pushforward-sheaf}
we have $(h_U^\#)_p = (h_U)_p$. Thus we see that $u$
is equal to the composition of functors
$$
\mathcal{C} \xrightarrow{h}
\textit{PSh}(\mathcal{C}) \xrightarrow{{}^\#}
\Sh(\mathcal{C}) \xrightarrow{()_p}
\textit{Sets}
$$
Each of these functors is left exact, and hence we see $u$ satisfies (1).

\medskip\noindent
Conversely, suppose that $u$ satisfies (1) and (2).
In this case we immediately see that $u$ satisfies the first two
conditions of Definition \ref{definition-point}. And its
stalk functor is exact, because it is a left adjoint by
Lemma \ref{lemma-point-pushforward-sheaf} and it commutes
with finite limits by Lemma \ref{lemma-neighbourhoods-directed}.
\end{proof}

\end{document}
